\section{本章小结}
\label{sec:block-summary}

本章从一个朴素的想法出发：变量合在一起难处理，拆开之后每一部分都便宜。全章回答了开头提出的两个问题。

第一，只动一部分变量能否推进完整目标。可以，但要看更新的类型。精确块更新靠“当前点本身是候选”保证下降，BCGD和BPG靠块上界保证下降，这些都可以统一到BSUM的下降链：上界不低于目标，当前点相等，最小化上界不会让目标上升。块Newton法和APCG不在这条链上，前者要靠线搜索或信赖域，后者靠辅助序列。

下面的表把四类更新放在一起。

\begin{BookInlineTable}
\centering
\begin{tabular}{p{18mm}p{30mm}p{46mm}}
\toprule
更新 & 子问题 & 主要条件\\
\midrule
精确BCD、ALS & 条件子问题 & 最小点存在，选解规则；ALS要求因子满列秩或用伪逆\\
BCGD & 二次上界 & 块光滑常数$L_J$\\
BPG & 二次上界加$g_J$ & $g_J$按块可分，邻近映射可算\\
APCG & 外推点的块邻近模型 & $f$凸，块光滑，解存在，均匀抽块，精确块求解\\
\bottomrule
\end{tabular}
\end{BookInlineTable}

下降不等于解出问题。逐块最优、驻点、全局最优是三个不同的层次。光滑或可分时逐块最优能合成驻点，联合凸时驻点才是全局最优；不可分的非光滑项会让坐标下降停在非驻点，分别凸的矩阵分解可能停在鞍点。循环规则下，聚点是逐块最优点需要紧水平集和子问题解唯一；独立均匀抽块的BCGD有有限步的梯度界，强凸时有线性收缩，都是对期望而言。

选块规则和块内更新可以自由组合，但保证要重新核对。随机选的是更新位置，不需要也不应当做逆概率校正；贪心规则的便宜程度取决于偏梯度能否增量维护；并行更新要用覆盖同时变化的联合上界，步长通常按同时更新的块数缩小。

第二，省下的计算会不会被更多的迭代抵消。这要落到总工作量上比较。维护残差后，一轮坐标更新和一次全梯度成本相同，这才让坐标方法在强耦合较少、块常数远小于全局常数时取得优势；块大小、选块、内层精度、外推和同步的成本都要算进去。

本章的偏梯度都是精确的。当目标是大量项的和、每次只读其中一部分时，方向本身带有误差，这是下一部分研究的问题。
